\documentclass[11pt,a4paper]{article}
\usepackage[utf8]{inputenc}
\usepackage[T1]{fontenc}
\usepackage{lmodern,amsmath,amssymb,amsthm,mathtools,booktabs,enumitem,microtype}
\usepackage[margin=1in]{geometry}
\usepackage{natbib}
\usepackage{xr-hyper}
\usepackage{hyperref}
\externaldocument[GE-]{01_law_of_graphic_equalisation_theory_referenced}
\externaldocument[MS-]{01b_graphic_equalisation_mind_and_self_grounding}
\externaldocument[QC-]{01f_graphic_equalisation_qualia_and_consciousness_grounding}
\externaldocument[NM-]{03_law_of_narrative_momentum_theory_referenced}

\newtheorem{definition}{Definition}[section]
\newtheorem{proposition}[definition]{Proposition}
\newtheorem{hypothesis}[definition]{Hypothesis}
\newtheorem{corollary}[definition]{Corollary}
\theoremstyle{remark}
\newtheorem{remark}[definition]{Remark}

\newcommand{\Frame}{\mathcal F}
\newcommand{\Supp}{\mathcal S}
\newcommand{\Risk}{\mathcal R}
\newcommand{\Cost}{\mathcal C}

\title{The Law of Graphic Equalisation --- Total Perspective Vortex Grounding\\\large Self-Model Dependence and the Cost of Cross-Frame Transition}
\author{Matthew Parslow\\Independent Researcher}
\date{\textbf{Draft \texorpdfstring{\textemdash}{---}\ \PaperDraft}\\[2pt]\PaperDisplayDate \texorpdfstring{\textbullet}{\textperiodcentered}\ document revision \PaperRevision}

\usepackage{fancyhdr}
\pagestyle{fancy}
\fancyhf{}
\fancyfoot[C]{\thepage}
\fancyfoot[R]{\scriptsize\PaperFooterStamp}
\renewcommand{\headrulewidth}{0pt}
\fancypagestyle{plain}{\fancyhf{}\fancyfoot[C]{\thepage}%
  \fancyfoot[R]{\scriptsize\PaperFooterStamp}\renewcommand{\headrulewidth}{0pt}}
% GENERATED by tools/build_paper_metadata.py from paper_metadata.json.
% DO NOT EDIT: edit paper_metadata.json, then regenerate this file.
\providecommand{\PaperDraft}{draft-2026-09-06.1}
\providecommand{\PaperDisplayDate}{6 September 2026}
\expandafter\def\csname PaperRevision@00_is\endcsname{129}
\expandafter\def\csname PaperRevision@01_law_of_graphic_equalisation_theory_referenced\endcsname{294}
\expandafter\def\csname PaperRevision@01a_graphic_equalisation_physics_grounding\endcsname{47}
\expandafter\def\csname PaperRevision@01b_graphic_equalisation_mind_and_self_grounding\endcsname{47}
\expandafter\def\csname PaperRevision@01c_graphic_equalisation_cosmology_grounding\endcsname{11}
\expandafter\def\csname PaperRevision@01d_graphic_equalisation_born_recovery\endcsname{6}
\expandafter\def\csname PaperRevision@01e_graphic_equalisation_communication_grounding\endcsname{6}
\expandafter\def\csname PaperRevision@01f_graphic_equalisation_qualia_and_consciousness_grounding\endcsname{3}
\expandafter\def\csname PaperRevision@01g_graphic_equalisation_total_perspective_vortex_grounding\endcsname{4}
\expandafter\def\csname PaperRevision@01h_graphic_equalisation_information_gravity_quantum_reconciliation\endcsname{5}
\expandafter\def\csname PaperRevision@01i_graphic_equalisation_passive_continuation_momentum_gravity\endcsname{12}
\expandafter\def\csname PaperRevision@01s_graphic_equalisation_statement\endcsname{294}
\expandafter\def\csname PaperRevision@02_seldon_manifolds_theory_referenced\endcsname{31}
\expandafter\def\csname PaperRevision@02a_seldon_manifolds_psychological_grounding\endcsname{15}
\expandafter\def\csname PaperRevision@02b_seldon_manifolds_social_and_psychological_application\endcsname{7}
\expandafter\def\csname PaperRevision@02c_seldon_manifolds_decision_making\endcsname{11}
\expandafter\def\csname PaperRevision@03_law_of_narrative_momentum_theory_referenced\endcsname{32}
\expandafter\def\csname PaperRevision@04_narrativium_theory_referenced\endcsname{23}
\expandafter\def\csname PaperRevision@05_frames_relation_and_the_empty_limit\endcsname{11}
\expandafter\def\csname PaperRevision@06_assessing_an_account_of_this_kind\endcsname{29}
\expandafter\def\csname PaperRevision@07_causal_attitude_networks_on_the_substrate\endcsname{4}
\expandafter\def\csname PaperRevision@08_certified_answers_as_an_audit\endcsname{4}
\expandafter\def\csname PaperRevision@09_adaptability\endcsname{94}
\expandafter\def\csname PaperRevision@09a_adaptability_implementation\endcsname{7}
\expandafter\def\csname PaperRevision@10_applied_graphic_equalisation\endcsname{7}
\expandafter\def\csname PaperRevision@10a_red_queen_race_civilisational_rate_of_change_llms_biological_timescales_anthropocene_transition\endcsname{13}
\expandafter\def\csname PaperRevision@11_relational_information_ethics\endcsname{16}
\expandafter\def\csname PaperRevision@A1_constructive_graphic_equalisation_implementation_referenced\endcsname{8}
\expandafter\def\csname PaperGenerated@01s_graphic_equalisation_statement\endcsname{1}
\edef\PaperMetadataStem{\jobname}
\expandafter\ifx\csname PaperRevision@\PaperMetadataStem\endcsname\relax
  \PackageError{paper_metadata}{No metadata for \PaperMetadataStem}{Add the paper stem to paper_metadata.json.}
  \def\PaperRevision{?}
\else
  \expandafter\let\expandafter\PaperRevision\csname PaperRevision@\PaperMetadataStem\endcsname
\fi
\newcommand{\PaperDateLine}{\textbf{Draft \texorpdfstring{\textemdash}{---}\ \PaperDraft}\\[2pt]\PaperDisplayDate\ \texorpdfstring{\textbullet}{\textperiodcentered}\ document revision \PaperRevision}
\newcommand{\PaperStatementDateLine}{\textbf{Draft \texorpdfstring{\textemdash}{---}\ \PaperDraft}\\[2pt]\PaperDisplayDate\ \texorpdfstring{\textbullet}{\textperiodcentered}\ statement}
\expandafter\ifx\csname PaperGenerated@\PaperMetadataStem\endcsname\relax
  \newcommand{\PaperFooterStamp}{\textsc{\PaperDraft\,r\PaperRevision}}
\else
  \newcommand{\PaperFooterStamp}{\textsc{\PaperDraft\ / statement / from long r\PaperRevision}}
\fi
\newcommand{\PaperNoticeStamp}{\textbf{DRAFT --- \PaperDraft\ r\PaperRevision} (\PaperDisplayDate).}


\begin{document}
\maketitle
\begin{center}
\fbox{\parbox{0.92\textwidth}{\small
\PaperNoticeStamp\ This document is the \emph{01g} grounding companion to
\emph{The Law of Graphic Equalisation --- Theory}, \emph{Mind and Self
Grounding}, \emph{Qualia and Consciousness Grounding}, and \emph{The Law of
Narrative Momentum}. It defines the Total Perspective Vortex as a recursive
failure of a demand for total perspective. Its positive account is bounded:
frames construct models, bridges, and costly transitions under support
conditions. It neither makes a clinical classification nor treats a model of
another frame as possession of that frame's experience.}}
\end{center}
\vspace{1em}

\begin{abstract}
The phrase ``Total Perspective Vortex'' can sound as though it names an
extraordinary mental state. This paper uses it instead as a defined formal term:
the recursive failure produced by demanding a complete, frame-free total
perspective from a finite frame. Graphic Equalisation states the situated-view premise (Axiom~\ref{GE-ax:no-view}), and Mind and Self Grounding develops self-model limits (Proposition~\ref{MS-prop:selfmodel}). This paper does not reopen either result. It
tracks their joint consequence for moving between perspectives. A frame $F$
organises an item $x$ as $P_F(x)$, carries a self-model $M_F(F)$, and can infer a
model $M_F(K)$ of another frame $K$ without thereby instantiating $K$'s
perspective. A mediating frame $W$ can construct a bounded bridge $B^W_{F,K}$
and a representational transition $\tau^W_{F\to K}$. Such a transition is not a
leap outside perspective: it is a new situated construction, assessed relative
to support $S_W$.

Narrative Momentum supplies the cost form. Reusing carried organisation is the
free continuation; altering it requires a specified intervention. Accordingly,
this paper defines conditional transition cost
$C_\tau(\tau\mid S_W)$ and a conditional self-model risk ordering. The latter
increases, within a specified grounding and with relevant support held fixed,
when a transition changes an organisation on which the frame's self-model is
more dependent, central or coupled, less redundant, and less supported. It is a
conceptual conditional quantity, not a clinical score, a universal numerical
law, evidence that the transition is true, or a credential for any experience.
The paper ends by specifying calibration tests and falsifiers for the bridge and
cost claims.
\end{abstract}

The literary device presents a total creation with a ``you are here'' marker;
its survival case is a universe constructed around the viewer. The analogy
motivates the distinction between totalisation and a view situated in a frame;
it supplies no clinical evidence or risk scale.

\section{Scope, dependency, and terminology guard}\label{sec:scope}
This is a focused Graphic Equalisation companion. It inherits the relational
floor, the running graph, frame-relative grounding, and refusal of a view from
nowhere from \emph{Graphic Equalisation}; it inherits the structural
incompleteness of self-models from \emph{Mind and Self Grounding}; and it
inherits the reuse-versus-alteration cost form from \emph{Narrative Momentum}.
It adds neither a consciousness theory nor a general proof that self-models are
incomplete. Its narrower question is what a frame can construct when it tries to
relate its perspective to another perspective without smuggling in a total view.

\paragraph{Terminological guard.} \emph{Total Perspective Vortex} (TPV) is a
term adapted from Douglas Adams's fictional device \citep{adams1980}; its formal definition here is this paper's own. It is not a clinical or established scientific term. It names a failed demand and its recursive repair attempt, not
a diagnosis, a disorder, a special capacity, or a credential. In particular, this
paper does not claim that derealisation, disorientation, distress, or any other
reported condition is required for, evidences, or authenticates the construction
described here. Questions of clinical meaning, care, prevalence, and measurement
are outside its scope.

Three distinctions keep the claim bounded. First, $M_F(K)$ is an inferred model
of $K$ \emph{in $F$}; it is not direct possession of $K$'s experience, a
frame-free comparison class, or proof that $F$ has reproduced every relation
operative in $K$. Second, a transition cost is a cost of alteration under named
conditions, not a measure of truth, value, courage, severity, or reality.
Third, the formal order introduced below is conditional on a chosen grounding;
it is not an empirical numerical law until a grounding supplies observables,
units, and a measurement procedure.

The nearest established vocabulary is adjacent rather than identical. Predictive
coding models treat perceptual processing as hierarchically organised prediction
and residual error \citep{raoballard1999,friston2010}; Metzinger's account makes
self-model transparency a central constraint \citep{metzinger2003,metzinger2020};
and metacognition research studies measurable divergence between confidence and
performance \citep{fleming2024}. Those accounts concern particular theoretical or
empirical machinery. The contribution here is the arrangement of frame-relative
perspective, bridge construction, and conditional alteration cost. It makes no
claim that any cited account establishes the present formalism.

\section{Frames, aperture, and perspective}\label{sec:frames}
A frame is a locus at which distinctions are made. Its aperture is the set and
grain of distinctions it can currently organise; it is not a hidden container of
all distinctions it does not presently hold. Let $x$ denote an item, event, or
relation as individuated at a containing level.

\begin{definition}[Frame-relative perspectival organisation]
For a frame $F$, \emph{frame-relative perspectival organisation} is
\[
 P_F(x)=\operatorname{Org}_F(x;A_F,H_F,\Gamma_F),
\]
where $A_F$ is F's operative aperture, $H_F$ its carried, currently operative
history-bearing structure, and $\Gamma_F$ its admissibility grammar. $P_F(x)$ is
not an item copied unchanged into $F$; it is the organisation $F$ can make and
use of $x$ under those conditions.
\end{definition}

This notation does not turn perspective into an inner picture. It names an
operative organisation: distinctions that are available, related, and
consequential to F's continuation. The same $x$ can support different $P_F(x)$
and $P_K(x)$ without either difference licensing a frame-free residual. A
containing frame can compare relations that neither lower frame represents, but
that comparison is $P_W(\cdot)$ and therefore carries W's aperture and grammar.

\begin{remark}[Aperture is a boundary, not a defect]
A finite aperture neither makes a frame irrational nor prevents it from learning,
modelling, or correcting. It fixes the condition under which a perspective is a
perspective: F's organisation is made from a location with a finite operative
vocabulary. An absent distinction is not thereby a value of zero, and an
unavailable relation is not direct evidence that no frame can hold it.
\end{remark}

\section{Carried trajectory, self-model, and model of another}\label{sec:models}
A frame carries a trajectory through its own running organisation. To use that
trajectory as an object of reasoning, it constructs a representational model;
that model is not identical with the carried course. This is the inherited
self-model distinction, used here only as a dependency.

\begin{definition}[Self-model and other-model]
$M_F(F)$ is the model of F that F can construct, retain, and use under its
operative aperture. For another frame $K$, $M_F(K)$ is F's inferred model of K:
the organisation F uses to predict, interpret, coordinate with, or otherwise
reason about K.
\end{definition}

The subscripts are load-bearing. $M_F(K)$ may incorporate reports, observed
consequences, shared records, and a model of how K organises an item. It remains
F's model. Replacing $M_F(K)$ with $P_K(x)$ erases the relation that needs
explaining: F has constructed a representation \emph{about} K, not stepped into
a view from nowhere or acquired K's direct experience. Conversely, treating the
model as worthless because it is not identity is another error: a model can be
calibrated, useful, and revisable without being total.

\begin{definition}[Carried trajectory]
Let $T_F=\langle P_{F,t}(x)\rangle_{t\leq n}$ denote the selectively retained,
currently operative succession of F's perspectival organisation. The carried
trajectory is the part of $T_F$ represented in present support and available to
shape F's next organisation.
\end{definition}

The qualification ``selectively retained'' is essential. Narrative Momentum does
not treat all past states as equally carried; it treats the history-bearing
structure still operative in the present as what persists. Hysteresis is a close
example: present response can retain selected branch or extremum information
without retaining a total history \citep{mayergoyz1991,mayergoyz2003}. In an
inferential setting, belief revision likewise studies how a change can preserve
as much of an existing set as the specified revision permits \citep{agm1985}.
Neither literature proves the present account; both identify the reuse and
alteration structure that the account requires.

\section{The Total Perspective Vortex and non-totalisation}\label{sec:vortex}
The demand that begins the vortex is: ``construct one perspective that is
complete, frame-free, and total over all perspectives.'' A finite frame can
respond by adding a model of itself, then a model of that model, then a
containing observer of both. The recursion changes the represented object but
never supplies the demanded location.

\begin{definition}[Total Perspective Vortex]
The \emph{Total Perspective Vortex} is the recursive failure of a frame's demand
for a complete, frame-free total perspective: every proposed totalising model is
constructed in a finite frame, so its conditions of construction are relocated
rather than eliminated; adding a model of those conditions restarts the same
demand at the containing organisation.
\end{definition}

\begin{proposition}[Non-totalisation]\label{prop:nontotalisation}
No finite frame obtains a complete, frame-free total perspective. A containing
observer relocates the aperture and self-model dependency from F to the
containing frame; it does not solve them.
\end{proposition}

This is not a second proof of generic self-model incompleteness. It is the direct
consequence of applying the inherited frame condition to the stronger demand for
totality. If $W$ represents F and K, W can construct relations unavailable to
either. But its representation is $P_W(\cdot)$, conditioned by $A_W$, $H_W$, and
$\Gamma_W$. Naming a still larger $V$ repeats the relocation. The result does
not forbid increasingly adequate models or multi-frame coordination. It forbids
calling any finite, situated construction the completion of a frame-free total
perspective.

\begin{corollary}[No vortex is required for ordinary modelling]
F need not demand totalisation in order to model K, revise a self-model, or
construct a useful bridge. The vortex is a failure mode of an impossible demand,
not the price of perspective-taking itself.
\end{corollary}

\section{Jumping off the tower: the cross-frame bridge}\label{sec:bridge}
``Jumping off the tower'' names the practical formal move: stop trying to
complete an infinite tower of observers and instead specify the finite bridge a
containing frame can actually construct. This is not an escape from frames. It
is a change from totalisation to accountable mediation.

\begin{definition}[Cross-frame bridge]
Given frames F and K, a \emph{cross-frame bridge} in mediating frame W is
\[
 B^W_{F,K}=\langle E^W_{F,K},\ \Pi^W_{F,K},\ \Lambda^W_{F,K}\rangle,
\]
where $E^W_{F,K}=\langle E^W_F,E^W_K\rangle$ is the evidence, record, or
contact relation W admits concerning F and K; $\Pi^W_{F,K}$ is W's
correspondence or translation procedure; and $\Lambda^W_{F,K}$ states the
bridge's loss, grain, and conditions of use. From its admitted evidence, W can
construct $M_W(F)$, $M_W(K)$, $\widehat P^W_F(x)$, and
$\widehat P^W_K(x)$. These W-indices are load-bearing: they are bounded
constructions about existing frames, not F- or K-local organisation imported
without mediation.
\end{definition}

A bridge can use shared language, behavioural regularities, measurements,
reports, jointly constrained tasks, or formally specified mappings. Its quality
is not guaranteed by its existence. The bridge must state what it preserves, what
it leaves unrepresented, and which counterevidence would require revision. The
model-of-another claim is therefore weaker and more useful than possession: it
makes coordination and correction possible while retaining the difference
between a construction and the perspective it models.

The derivation posits F and K as existing frames; it does not grant W exclusive,
complete, or direct understanding of either. $P_F(x)$ remains the organisation
operative in F, but it is not an input datum available to a W-mediated procedure.
W works through $E^W_F$ and the constructions it can make from that admitted
relation. Another containing frame can construct a different or finer relation
to F without any one representation exhausting F.

\begin{definition}[Representational transition]
A \emph{representational transition} through W is a procedure
\[
 \tau^W_{F\to K}:\bigl(E^W_F,M_W(F),B^W_{F,K}\bigr)
 \longmapsto \widehat P^W_K(x),
\]
that produces a W-indexed, bridge-limited representation $\widehat P^W_K(x)$ of
how K might organise x. The arrow denotes constructed transformation, not
transport of F into K or acquisition of K's experience.
\end{definition}

The transition is deliberately asymmetric. $\tau^W_{F\to K}$ begins with W's
admitted relation to F and ends in a representation W can hold about K. A reverse
construction has its own bridge, losses, and conditions. Equality, invertibility,
or mutual transparency cannot be inferred from notation alone.

\section{Transition cost and self-model risk}\label{sec:cost-risk}
The bridge answers what can be constructed; it does not say what must be altered
to use the construction. That question uses Narrative Momentum's conditional alteration-cost comparison (Hypothesis~\ref{NM-hyp:narrative-momentum}).
A continuation that reuses the relevant carried structure is free with respect to
that alteration; a continuation that requires the structure to be rewritten,
bypassed, or made inoperative pays the specified alteration cost. Path-dependence
and increasing-returns research provides familiar local instances in which
historically selected structures become harder to abandon \citep{arthur1989,david1985,pierson2000}.

\begin{definition}[Conditional transition cost]
A \emph{cost grounding} $\mathfrak C$ specifies a transition comparison class,
a cost domain $(\mathbb C_{\mathfrak C},\preceq_{\mathfrak C})$, its units or
ordinal rule, and the support configuration $S_W$ under which comparisons are
made. The \emph{conditional transition cost} is the value
\[
 C^{\mathfrak C}_\tau\bigl(\tau^W_{F\to K}\mid S_W\bigr)
 \in\mathbb C_{\mathfrak C},
\]
which $\mathfrak C$ assigns to constructing and maintaining the transition while
preserving its stated conditions. It is undefined when no comparison class, cost
order, or usable support condition has been supplied.
\end{definition}

For a chosen transition class, let $H^\star$ be the carried organisation relevant
to that class. Its reuse-versus-alteration form is
\[
 C^{\mathfrak C}_\tau(\tau\mid S_W)=
 C^{\mathfrak C}_{\mathrm{reuse}}(\tau\mid S_W)
 \oplus \Delta^{\mathfrak C}_{\mathrm{alter}}(H^\star,\tau\mid S_W),
 \qquad \Delta^{\mathfrak C}_{\mathrm{alter}}\succeq_{\mathfrak C}0.
\]
Here $\Delta^{\mathfrak C}_{\mathrm{alter}}=0$ when $\tau$ reuses $H^\star$.
Where $\tau$ requires $H^\star$ to be altered, bypassed, or made inoperative,
$\Delta^{\mathfrak C}_{\mathrm{alter}}$ is the specified alteration cost. This
equality applies only when $\mathfrak C$ supplies an additive cost currency.
Where it supplies only an ordinal comparison, the invariant claim is narrower:
an otherwise matched transition that reuses $H^\star$ ranks no higher in
alteration burden than one that requires its alteration. Neither form assumes a
universal, monetary, conscious, or shared unit. In the free-continuation case,
zero alteration increment in an additive implementation does not imply that
total activity is costless in a containing frame.

\begin{remark}[Reuse does not make a transition truth-preserving]
Reuse identifies the carried organisation that avoids further alteration cost in
the selected class. It does not establish that the organisation is adequate,
true, or normatively preferable. A bridge may reuse a badly calibrated model and
still require revision on independent evidence.
\end{remark}

\begin{definition}[Conditional self-model risk ordering]
A \emph{risk grounding} $\mathfrak R$ specifies a comparison class
$\mathcal T_{\mathfrak R}$, an ordered risk domain
$(\mathbb Q_{\mathfrak R},\preceq_{\mathfrak R})$, four ordered factor domains
$(\mathbb D_{\mathfrak R},\preceq_D)$,
$(\mathbb C_{\mathfrak R},\preceq_C)$,
$(\mathbb R_{\mathfrak R},\preceq_R)$, and
$(\mathbb S_{\mathfrak R},\preceq_S)$, and maps
\[
 d,c,r,s:\mathcal T_{\mathfrak R}
 \longrightarrow
 \mathbb D_{\mathfrak R},\mathbb C_{\mathfrak R},
 \mathbb R_{\mathfrak R},\mathbb S_{\mathfrak R},
\]
respectively measuring self-model dependence, centrality or coupling,
redundancy, and post-transition remaining support. Each factor order is oriented
so that larger elements mean more of the named factor. The initial support
configuration is $S_W^0$; $s_\tau$ is what remains after the particular
transition and is therefore not held fixed by conditioning on $S_W^0$. The risk
map is a grounding-supplied map, with no monotonicity assumed by its type
\[
 \Phi_{\mathfrak R}:
 \mathbb D_{\mathfrak R}\times\mathbb C_{\mathfrak R}
 \times\mathbb R_{\mathfrak R}^{\mathrm{op}}
 \times\mathbb S_{\mathfrak R}^{\mathrm{op}}
 \longrightarrow\mathbb Q_{\mathfrak R},
\]
and defines
\[
 \Risk^{W,\mathfrak R}_F(\tau\mid S_W^0)=
 \Phi_{\mathfrak R}(d_\tau,c_\tau,r_\tau,s_\tau).
\]
This orders vulnerability to loss of calibration, continuity, or bridge use under
the selected grounding; it is not a clinical score or a universal numerical law.
\end{definition}

\begin{hypothesis}[Conditional dependence risk]\label{hyp:conditional-risk}
For $\tau_1,\tau_2\in\mathcal T_{\mathfrak R}$ with the same target,
bridge-quality criterion, initial support $S_W^0$, and measurement grounding,
if $d_{\tau_1}\succeq_D d_{\tau_2}$,
$c_{\tau_1}\succeq_C c_{\tau_2}$,
$r_{\tau_1}\preceq_R r_{\tau_2}$, and
$s_{\tau_1}\preceq_S s_{\tau_2}$, with at least one strict inequality, then
\[
 \Risk^{W,\mathfrak R}_F(\tau_2\mid S_W^0)
 \prec_{\mathfrak R}
 \Risk^{W,\mathfrak R}_F(\tau_1\mid S_W^0).
\]
An intervention may add initial support or redundancy, thereby defining a new
comparison configuration; it cannot be described as holding $S_W^0$ fixed.
Neither a higher risk order nor a higher transition cost establishes the truth,
value, or clinical meaning of the transition.
\end{hypothesis}

The conditionals do the work. The order applies only after a risk grounding
supplies its comparison class, observables, and ordering rule. A false model can
be expensive to alter and a true correction can be cheap when support and
redundancy are strong. Conversely, a low cost neither proves safety nor truth.
The claim is that the named configuration determines why formally similar
transitions can have different alteration burdens, and that this claim can break
at a specified measurement boundary.

\section{Failure modes, calibration, tests, and falsifiers}\label{sec:tests}
The account fails if its limits are silently removed. Four recurrent mistakes
provide direct calibration checks.
\begin{enumerate}[leftmargin=*]
\item \emph{Totalisation error:} treat a containing model as a frame-free view.
The repair is to index the construction to W and state $A_W$, $\Gamma_W$, and
its bridge conditions.
\item \emph{Possession error:} treat $M_W(K)$ or $\widehat P^W_K(x)$ as K's
direct experience. The repair is to preserve the model, bridge, and target
indices and list unrepresented relations.
\item \emph{Cost-as-truth error:} infer that difficulty authenticates the
transition. The repair is an independent bridge calibration: compare its
predictions or coordination outcomes with admitted evidence, not with its cost.
\item \emph{Clinical reification error:} promote the formal variables into a
diagnosis, a severity ranking, or a required phenomenology. The repair is to
keep them as conditional structural quantities until an appropriate empirical
and ethical measurement programme independently licenses more.
\end{enumerate}

\paragraph{Calibration protocol.} For a specified application, first declare a
cost grounding $\mathfrak C$ and risk grounding $\mathfrak R$: their comparison
classes, units or ordering rules, initial support $S_W^0$, and independently
observable measures for $d_\tau,c_\tau,r_\tau,s_\tau$. Then construct at least
two bridges with their losses stated. A bridge earns use only insofar as it
improves specified prediction, coordination, compression, or error detection
relative to an available comparison; it is not validated by eloquence or by an
assertion of empathy. Compare transitions within the declared class. An
intervention that adds redundancy or support creates a new initial configuration;
the hypothesis predicts lower ordered risk only when the added relation
independently preserves or replaces the relevant support, not when it is merely
additional description.

\begin{proposition}[Falsifiers]\label{prop:falsifiers}
The framework's claims are exposed at three levels. The non-totalisation result
fails if a finite construction can supply a complete, frame-free perspective
while retaining no unrepresented construction conditions. The bridge account
fails if a purportedly unmediated transition can be shown to contain neither
admitted evidence, translation, nor loss conditions while still relating
non-identical apertures. The conditional-risk hypothesis fails in a grounded measurement programme if,
within its declared comparison class and common initial support $S_W^0$,
independently measured increases in dependence or coupling, decreases in
redundancy, and decreases in post-transition remaining support fail to order
risk as specified across the relevant cases.
\end{proposition}

The last exposure is deliberately not a universal prediction. A negative result
in one implementation can locate a bad operationalisation of $d,c,r,s$, a
missing support variable, or a failure of the hypothesised ordering in that
class. It cannot be repaired by declaring cost to be truth or by relabelling an
ordinary disagreement as a vortex.

\section{Prior art, relation to 01b/01f/03, and conclusion}\label{sec:conclusion}
This paper's relation to its neighbouring documents is division of labour.
\emph{Mind and Self Grounding} establishes that a self-model is constructed
under a finite aperture and distinguishes calibration from completeness; 01g
uses that result as the source of self-model dependence but does not repeat its
generic proof. \emph{Qualia and Consciousness Grounding} concerns whether
determinate phenomenal organisation requires a second ontological kind; 01g
does not reopen that question and does not infer experience from bridge
construction. \emph{Narrative Momentum} supplies the present-structure and
reuse-versus-alteration form; 01g applies it specifically to cross-frame
representational transition rather than claiming a new law of history.

The adjacent research record supplies pieces, not a hidden proof. Predictive
coding supplies models of hierarchical inference and residual updating
\citep{raoballard1999,friston2010}. Metzinger develops a substantial account of
self-model transparency \citep{metzinger2003,metzinger2020}; Fleming reviews
measurable calibration in metacognitive confidence \citep{fleming2024};
hysteresis, path dependence, and AGM revision provide established analyses of
retained state, lock-in, and conservative change
\citep{mayergoyz1991,mayergoyz2003,arthur1989,david1985,pierson2000,agm1985}.
None establishes a complete perspective, and none is claimed to validate a
clinical use of TPV.

The demystification is exact. The Total Perspective Vortex is not a hidden realm
behind ordinary perspective-taking. It is what occurs when a finite frame demands
a view that would have to erase the very aperture and construction through which
it is made. The alternative is not resignation or solipsism. It is finite,
inspectable bridge construction: model another frame as another frame; state the
bridge and its losses; reuse what carries forward where that is appropriate; and
name the conditional cost of altering what the self-model depends on. A
containing observer can improve the bridge, but it relocates rather than abolishes
the conditions of perspective. That boundary is the result.

\begin{thebibliography}{99}

\bibitem[Adams(1980)]{adams1980}
Douglas Adams.
\newblock \emph{The Restaurant at the End of the Universe}.
\newblock Pan Books, 1980.

\bibitem[Alchourr\'on et~al.(1985)]{agm1985}
Carlos E. Alchourr\'on, Peter G\"ardenfors, and David Makinson.
\newblock On the logic of theory change: Partial meet contraction and revision functions.
\newblock \emph{Journal of Symbolic Logic}, 50(2):510--530, 1985.

\bibitem[Arthur(1989)]{arthur1989}
W. Brian Arthur.
\newblock Competing technologies, increasing returns, and lock-in by historical events.
\newblock \emph{The Economic Journal}, 99(394):116--131, 1989.

\bibitem[David(1985)]{david1985}
Paul A. David.
\newblock Clio and the economics of QWERTY.
\newblock \emph{American Economic Review}, 75(2):332--337, 1985.

\bibitem[Fleming(2024)]{fleming2024}
Stephen M. Fleming.
\newblock Metacognition and confidence: A review and synthesis.
\newblock \emph{Annual Review of Psychology}, 75(1):241--268, 2024.

\bibitem[Friston(2010)]{friston2010}
Karl Friston.
\newblock The free-energy principle: a unified brain theory?
\newblock \emph{Nature Reviews Neuroscience}, 11(2):127--138, 2010.

\bibitem[Mayergoyz(1991)]{mayergoyz1991}
Isaak D. Mayergoyz.
\newblock \emph{Mathematical Models of Hysteresis}.
\newblock Springer, New York, 1991.

\bibitem[Mayergoyz(2003)]{mayergoyz2003}
Isaak D. Mayergoyz.
\newblock \emph{Mathematical Models of Hysteresis and Their Applications}.
\newblock Academic Press, second edition, 2003.

\bibitem[Metzinger(2003)]{metzinger2003}
Thomas Metzinger.
\newblock \emph{Being No One: The Self-Model Theory of Subjectivity}.
\newblock MIT Press, Cambridge, MA, 2003.

\bibitem[Metzinger(2020)]{metzinger2020}
Thomas Metzinger.
\newblock Self-modeling epistemic spaces and the contraction principle.
\newblock \emph{Cognitive Neuropsychology}, 37(3--4):197--201, 2020.

\bibitem[Pierson(2000)]{pierson2000}
Paul Pierson.
\newblock Increasing returns, path dependence, and the study of politics.
\newblock \emph{American Political Science Review}, 94(2):251--267, 2000.

\bibitem[Rao and Ballard(1999)]{raoballard1999}
Rajesh P. N. Rao and Dana H. Ballard.
\newblock Predictive coding in the visual cortex: a functional interpretation of some extra-classical receptive-field effects.
\newblock \emph{Nature Neuroscience}, 2:79--87, 1999.

\end{thebibliography}

\end{document}
